\begin{tikzpicture}[font=\small]
  % ---------- SRTBOT ----------
  \begin{scope}
    \node[anchor=south, font=\small] at (6.0,4.6){\textbf{SRTBOT}：6.006 给动态规划定的六步模板};

    \foreach \rw/\letter/\name/\what in {
      0/{\textbf{S}}/{Subproblem}/{定义子问题 —— \textbf{最难也最关键}的一步，其余五步几乎自动},
      1/{\textbf{R}}/{Relate}/{写出子问题之间的递推关系},
      2/{\textbf{T}}/{Topological order}/{论证子问题之间无环，从而可以按序求解},
      3/{\textbf{B}}/{Base case}/{给出边界情形},
      4/{\textbf{O}}/{Original problem}/{说明原问题由哪些子问题拼出},
      5/{\textbf{T}}/{Time}/{子问题个数 $\times$ 单个子问题的代价}}{
      \node[blk, minimum width=0.9cm, minimum height=0.62cm, fill=algblue!12, draw=algblue]
        at (0,3.8-\rw*0.78) {\letter};
      \node[anchor=west, font=\small] at (0.7,3.8-\rw*0.78) {\name};
      \node[note, anchor=west, align=left] at (3.9,3.8-\rw*0.78) {\what};
    }
  \end{scope}

  % ---------- subproblem DAG ----------
  \begin{scope}[yshift=-3.8cm, xshift=0.6cm]
    \node[anchor=south, font=\small] at (4.4,2.4){子问题 DAG：DP 就是在这张图上做\textbf{记忆化的 DFS}};

    \node[vtx, minimum size=8mm] (f5) at (0,1.2) {$F_5$};
    \node[vtx, minimum size=8mm] (f4) at (2.0,1.8) {$F_4$};
    \node[vtx, minimum size=8mm] (f3) at (2.0,0.6) {$F_3$};
    \node[vtx, minimum size=8mm] (f2) at (4.0,1.2) {$F_2$};
    \node[done, minimum size=8mm] (f1) at (6.0,1.8) {$F_1$};
    \node[done, minimum size=8mm] (f0) at (6.0,0.6) {$F_0$};

    \draw[dedge] (f5) -- (f4); \draw[dedge] (f5) -- (f3);
    \draw[dedge] (f4) -- (f3); \draw[dedge] (f4) -- (f2);
    \draw[dedge] (f3) -- (f2); \draw[dedge] (f3) -- (f1);
    \draw[dedge] (f2) -- (f1); \draw[dedge] (f2) -- (f0);

    \node[note, anchor=west, align=left] at (7.2,1.2)
      {\textcolor{algteal}{绿}：base case\\[4pt]
       朴素递归会把 $F_3$、$F_2$ 各算很多遍\\
       $\Rightarrow$ 指数时间。\\[4pt]
       记忆化让每个结点\textbf{只算一次}\\
       $\Rightarrow$ 时间 $=$ 结点数 $\times$ 每点出边数};

    \node[note, anchor=north, align=center] at (4.2,-0.2)
      {「子问题之间无环」这一条不是形式要求，而是\textbf{可解性的前提} ——
       有环就意味着子问题互相依赖，没有求解顺序。};
  \end{scope}

  % ---------- top-down vs bottom-up ----------
  \begin{scope}[yshift=-8.2cm, xshift=0.6cm]
    \node[blk, minimum width=4.4cm, minimum height=1.5cm, align=left, anchor=north west,
          fill=algteal!7, draw=algteal] at (0,1.0)
      {\textbf{自顶向下（记忆化）}\\[2pt]
       \footnotesize 照着递推直接写递归 + 缓存\\
       \footnotesize 只算到真正需要的子问题\\
       \footnotesize 递归深度可能爆栈};
    \node[blk, minimum width=4.4cm, minimum height=1.5cm, align=left, anchor=north west,
          fill=algblue!7, draw=algblue] at (5.2,1.0)
      {\textbf{自底向上（填表）}\\[2pt]
       \footnotesize 按拓扑序循环填表\\
       \footnotesize 无递归开销，常数更小\\
       \footnotesize 便于做滚动数组省空间};
    \node[note, anchor=north, align=center] at (4.8,-0.8)
      {两者时间复杂度相同，先用记忆化把递推写对，需要时再改写成填表。};
  \end{scope}
\end{tikzpicture}
